%! Quadratic Functions Revisited — Unit 2, Lesson 2.1 — Exit Ticket (Answer Key)
\documentclass[10pt]{article}
\usepackage{precalculus-article}
\usepackage{precalculus-key}   % pulls in -boxes; do NOT also load -boxes
\newcommand{\TallMath}[1]{$\displaystyle #1\rule[-1.4em]{0pt}{3.2em}$}

\begin{document}

\pageheader{Unit 2, Lesson 2.1}{Exit Ticket --- Answer Key}

\vspace{0.12in}

\begin{enumerate}[itemsep=12pt]

\item Let $f(x) = x^2 - 2x$. Find the average rate of change over $[1, 4]$.

\begin{work}
  f(1) &= -1 \\
  f(4) &= 8 \\
  \text{AROC} &= \frac{8 - (-1)}{4 - 1} \\
              &= 3
\end{work}

AROC $= $ \ans{3}

\item Two functions are given at inputs that step by $1$. Their first
      differences are:

      Function A: $6,\ 6,\ 6,\ 6$ \qquad Function B: $2,\ 6,\ 10,\ 14$

\begin{enumerate}[label=(\alph*), itemsep=4pt]
  \item Which one is linear? \ans{A} \quad Which one is quadratic?
        \ans{B}
  \item Second differences of the quadratic one: \ans{4} , \ans{4} ,
        \ans{4}
  \item In one sentence, justify your classification using those numbers.

        \vspace{0.05in}
        \ansline{A's rates of change never change, which is what linear means; B's rates change}
        \ansline{by a constant 4 every step, and constant second differences mean quadratic.}
\end{enumerate}

\item Let $f(x) = -3(x-2)^2 + 7$.

\begin{enumerate}[label=(\alph*), itemsep=4pt]
  \item Vertex: ( \ans{2} , \ans{7} ) \quad Axis of symmetry:
        $x = $ \ans{2}
  \item The graph is concave \ans{down}
\end{enumerate}

\end{enumerate}

\end{document}
